\documentclass[10pt,a4paper]{article}
\usepackage[margin=18mm]{geometry}
\usepackage[T1]{fontenc}
\usepackage[utf8]{inputenc}
\usepackage{lmodern,amsmath,amssymb,booktabs,graphicx}
\setlength{\parindent}{0pt}
\setlength{\parskip}{4pt}
\newcommand{\ext}{\mathrm{ext}}
\newcommand{\SstarRe}{0.0059232}
\newcommand{\SstarIm}{0.0548189}
\newcommand{\SstarAbs}{0.05514}
\newcommand{\SmaxNum}{103}
\newcommand{\AoneLo}{1.3333}
\newcommand{\AoneHi}{1.3337}
\newcommand{\AtwoLo}{0.4566}
\newcommand{\AtwoHi}{0.4567}
\newcommand{\AthreeLo}{-0.268}
\newcommand{\AthreeHi}{-0.265}
\newcommand{\StrongFitSentence}{Fitting $(m-3s/4)\sqrt s=a_1+a_2s^{-1/2}+a_3s^{-1}+\dots$ (5--7 terms, $s\ge10,15,20$) gives $a_1\in[1.33330,1.33369]$ with all coefficients free; with $a_1=4/3$ fixed, $a_2\in[0.45658,0.45669]$ (conditional prediction $0.456527$) and $a_3\in[-0.268,-0.265]$. Fixing also $a_2=0.456527$ (six further terms) fits as well as the free fits (rms $\le1\times10^{-8}$) and gives $a_3\in[-0.2625,-0.2623]$.}
\newcommand{\AthreeFixLo}{-0.262}
\newcommand{\AthreeFixHi}{-0.262}
\newcommand{\GlobalKappa}{0.3562}
\newcommand{\GlobalErr}{1.1e-05}
\newcommand{\GlobalFitSentence}{A ten-term approximant with the exact ends (\S4) reproduces all 126 numerical points to 1.1e-05 (relative).}

\begin{document}

{\Large\bfseries Extremal mass of equal-spin EGB black holes in $D=5$:\\[2pt]
weak series to $s^6$, its singularity, and the strong-coupling limit}\\[3pt]
{\small 16--23 October 2026. Continuation of \texttt{extremal-mass-report.pdf}.
$G_5=1$, action $R+\alpha\mathcal L_{\rm GB}$, $J_1=J_2=J$,
$M_\ext=\pi^{1/3}J^{2/3}m(s)$, $s=\pi^{2/3}\alpha/J^{2/3}$.}

\section*{Summary}
\begin{enumerate}
\item \textbf{Weak coupling, two more exact orders.} $d_5$ and $d_6$ (below) from the exact
exterior program; the horizon shifts $X_5,X_6$ and the first law are checked at each order.
\item \textbf{The near-horizon data are algebraic.} The extremal throat equations fix $g_H(s)$
exactly; their series reproduces the perturbative horizon through $\alpha^4$ and predicted
$X_5,X_6$ before the exterior computation finished.
\item \textbf{Radius of the weak series.} $g_H(s)$ has its nearest branch points at
$s_*=\SstarRe\pm\SstarIm\,i$, $|s_*|=\SstarAbs$, a square-root fold of the throat.
The signs of $d_1,\dots,d_6$ are those of $\mathrm{Re}\,s_*^{-n}$ (6 of 6).
Conjecture: $m(s)$ has the same singularity, so the weak series converges only for $|s|<\SstarAbs$.
\item \textbf{Strong coupling from pure-Lovelock universality.} Static $D=5$ EGB is the
``derivative-degenerate'' Lovelock theory of Dadhich--Pons--Prabhu (arXiv:1201.4994); its
small-hole core is a BTZ black hole with $\ell^2=4\alpha$, $1/(4G_3)=3\pi\alpha$. Using the
extremal-BTZ relation $\Delta M=S^2/(48\pi^3\alpha^2)$ with the \emph{exact} extremal entropy gives
\[ M_\ext=\frac{3\pi\alpha}{4}+\frac{4J}{3\sqrt\alpha}+\dots,\qquad
m(s)=\frac{3s}{4}+\frac43 s^{-1/2}+\dots \]
which reproduces the coefficient $4/3$ of the earlier matched-asymptotics candidate by an
independent route.
\item \textbf{Numerics at strong coupling.} Two-sided shooting ($s\le8.6$) and multiple shooting
with quadruple-precision near-horizon segments reach $s=\SmaxNum$ (the spectral branch stopped at
$5.24$). \StrongFitSentence{} Hence, numerically,
\[ m(s)=\frac{3s}{4}+\frac43\,s^{-1/2}+a_2\,s^{-1}+a_3\,s^{-3/2}+O(s^{-2}),\qquad
a_2=0.4566\pm0.0002,\quad a_3\simeq-0.26, \]
confirming $a_1=4/3$ and the conditional matched-asymptotics value $a_2=0.456527$
(fits with $a_2$ fixed to it are as good as free fits).
\item \textbf{No closed form.} $\zeta(3),\zeta(5),\zeta(3)^2,\pi^{2k}$ in $d_n$ exclude an algebraic
$m(s)$; the complex singularity at $|s|=0.055$ excludes a simple elementary form built from the
two ends. \GlobalFitSentence
\end{enumerate}

\section*{1. Weak coupling through $s^6$}
\[
m(s)=\tfrac32+s-\tfrac{236}{105}s^2+d_3s^3+d_4s^4+d_5s^5+d_6s^6+O(s^7),
\]
\begin{align*}
d_5&=\tfrac{447853140007115845664}{6100990520625}+\tfrac{657983488}{175}\zeta(3)
-\tfrac{3972366336}{625}\zeta(2)^2-\tfrac{1464139776}{25}\zeta(5)=423.920197004910\ldots\\
d_6&=-\tfrac{1245254609028150217918650165824}{88766575114501659375}
+\tfrac{15735548018688}{6125}\zeta(2)^3-\tfrac{57616520085504}{109375}\zeta(2)^2
+\tfrac{5569511424}{35}\zeta(3)^2\\
&\quad-\tfrac{94884252024832}{165375}\zeta(3)+\tfrac{2699634147328}{625}\zeta(5)
=-4737.494826989936\ldots
\end{align*}
($d_3,d_4$ as in the earlier report.) Method and conditions are those of
\texttt{code/pert.py}: exact hyperlogarithm quadratures, constants fixed by asymptotic flatness
and a double zero of $g^{rr}$; exact \emph{conditional} on the PSLQ identification (100 digits)
of the regularised horizon constants $C_w$ with multiple zeta values. At orders 5 and 6
the first-law residual $dM=2\Omega_H\,dJ$ is identically zero, and the horizon shifts
$X_5=292028416/729$, $X_6=58879115264/6561$ equal the throat predictions of \S2.
Run times 25.5\,min (order 5) and 1\,h\,38\,min (order 6).

\section*{2. Algebraic near-horizon data and the radius of convergence}
With $a=\alpha/g_H$ and $q=h_H/g_H$ the extremal throat satisfies
$P(a,q)=16a^2q^2(4-q)-8a(2q^2-7q+4)+q-2=0$ and
$8J^2/(\pi^2g_H^3)=q^2(4-q)[1+4a(4-3q)]/8$.
At fixed $J$ these give $g_H(s)$ as an algebraic function, with
$g_H=1+\tfrac83 s+\tfrac{16}9 s^2-\tfrac{10496}{81}s^3-512 s^4+\tfrac{9125888}{729}s^5+\dots$
(in units where $g_H(0)=1$). Two independent checks: (i) these coefficients equal the
horizon shifts of the exterior computation through $s^6$; (ii) the map $s(q)$ reproduces
the coupling of every point of the earlier spectral branch from its $q$ alone, to $10^{-6}$.

Eliminating $(g_H,q)$ from the throat equations and their Jacobian gives the candidate branch
points; Pad\'e approximants of 60 exact Taylor coefficients select those on the physical
sheet: $s_*=\SstarRe\pm\SstarIm\,i$. In the $q$-plane $s_*$ is a critical point of the explicit
map $s(q)$, at $q_c=2.16110\mp0.94332\,i$: a square-root fold. The Pad\'e poles beyond
$s_*$ lie on the line $\mathrm{Re}\,s\simeq0.006$, i.e.\ a cut running away from $s_*$.

\emph{Evidence that $m(s)$ shares it.} $\mathrm{sign}\,d_n=\mathrm{sign}\,\mathrm{Re}\,s_*^{-n}$
for $n=1,\dots,6$ ($\arg s_*=83.8^\circ$, so the pattern is not trivial); $|d_n||s_*|^n$
decreases slowly. A conformal map that opens the cuts from $s_*,\bar s_*$ turns the divergent
series into one oscillating around the numerical $m(s)$ out to $s\approx0.5$
(\texttt{series/conformal.py}); its slow convergence indicates further singularities near the
boundary of the mapped disk. None of this is a proof.

\section*{3. Strong coupling: pure-Lovelock universality}
For $D=5$ EGB the static master equation is $\psi+2\alpha\psi^2=\mu/r^4$, $V=1-\psi r^2$;
completing the square it is the $N=2$ ``derivative-degenerate'' polynomial
$F(\psi)=-a+(\psi-\beta)^2$ of arXiv:1201.4994 (their eq.\ 15) with $\beta=-1/(4\alpha)$.
For $r\ll\sqrt\alpha$ it behaves like pure Gauss--Bonnet in its critical odd dimension,
which (their \S IV) has BTZ-like, universal thermodynamics: with $\mu=r_h^2+2\alpha$,
\[ \Delta M=\tfrac{3\pi}{8}r_h^2,\quad S=6\pi^2\alpha\, r_h,\quad T=\frac{r_h}{8\pi\alpha}
\;\Longleftrightarrow\; \text{BTZ with }\ell^2=4\alpha,\ \tfrac1{4G_3}=3\pi\alpha . \]
The critical static solution $\mu=2\alpha$, $M=3\pi\alpha/4$, is the massless BTZ point.
For an extremal BTZ hole with the same $(\ell,G_3)$, $\Delta M=4G_3S^2/(4\pi^2\ell^2)=S^2/(48\pi^3\alpha^2)$.
The extremal entropy is known exactly from the throat; in $t=(2a)^{-1/2}$
\[ \frac{S_\ext^2}{48\pi^3\alpha^{3/2}J}=\frac43+\frac38t^2+O(t^3)
\quad\Longrightarrow\quad \Delta M_\ext=\frac{4J}{3\sqrt\alpha}\,(1+O(t^2)) \]
(\texttt{strong/btz\_universality.py}, symbolic). The step that is assumed, not derived, is
that the rotating extremal core obeys the static BTZ dictionary. The absence of an $O(t)$ term
means the next coefficient $a_2$ of $s^{-1}$ in $m(s)$ cannot come from the core.

\section*{4. Numerical solution at large $s$}
\textbf{Methods.} (i) $s\le5.24$: the spectral branch of \texttt{work/kkr-extremal}.
(ii) $2\le s\le8.6$: two-sided shooting (\texttt{strong/twoside.py}): outward from the horizon,
with its formal series (rank-one level blocks) plus the non-analytic throat mode $B\,x^\gamma$
($\gamma$ the exact throat exponent), inward from the large-$r$ series with free mass $\mu$ and
squashing amplitude $b_4$; $m=3\mu/4$. (iii) $s\ge9$: multiple shooting (\texttt{strong/ms.py},
\texttt{ms\_cont.py}). Single shooting fails for $s\gtrsim10$ because in the core the squashing
mode behaves as $e^{\pm4\sqrt{6}\alpha/r}$ (a factor $e^{76}$ at $\alpha=10$). The solver places
nodes at equal WKB phase in the core, integrates segments with $x\le0.5\rho$ in quadruple
precision ($\_\_$float128 Bulirsch--Stoer) from 60-digit horizon series, and continues in $\alpha$
across the resonances $\gamma=3,4$ ($\alpha=4.83,13.97$).

\textbf{Checks.} At the common point $s=2.98202$ the two-sided and spectral values differ by
3.4e-07; against a spline of the spectral branch the shooting/BVP
values deviate by at most 2.0e-05, the size of the spectral branch's own constraint
violations (median 3e-06, max 7e-05). Varying the shooting parameters moves $m$ by
at most 2.8e-08.
Multiple-shooting Newton residuals are below 3e-10. A degree-6 polynomial through the
18 multiple-shooting points with $9.16\le s\le10.3$, extrapolated to the BVP point $s=9.0037$,
agrees with it to 2e-09.

\begin{center}\small\begin{tabular}{rrrl}\toprule $s$ & $m(s)$ & $(m-3s/4)\sqrt s$ & method\\\midrule
2 & 2.4446986200 & 1.33600560 & shoot\\
2.4 & 2.6802239675 & 1.36363711 & shoot\\
2.98202 & 3.0410066937 & 1.38923766 & shoot\\
3.9 & 3.6394756194 & 1.41097629 & shoot\\
5.07 & 4.4346318566 & 1.42334984 & shoot\\
6.591 & 5.4998298751 & 1.42890434 & shoot\\
8.5683 & 6.9145757941 & 1.42948377 & shoot\\
9.00366 & 7.2290351282 & 1.42916958 & BVP\\
9.16046 & 7.3425005931 & 1.42903362 & MS\\
9.32 & 7.4580467159 & 1.42888461 & MS\\
9.52 & 7.6030391456 & 1.42868400 & MS\\
9.59986 & 7.6609741162 & 1.42859997 & MS\\
9.86 & 7.8498672728 & 1.42831221 & MS\\
10.08 & 8.0097946897 & 1.42805387 & MS\\
11.0299 & 8.7020175240 & 1.42681967 & MS\\
13.5121 & 10.5212075494 & 1.42316670 & MS\\
16.5529 & 12.7633325405 & 1.41866050 & MS\\
20.278 & 15.5223943689 & 1.41364795 & MS\\
24.8414 & 18.9136032758 & 1.40838925 & MS\\
32.562 & 24.6670488835 & 1.40131488 & MS\\
39.8898 & 30.1384093228 & 1.39612751 & MS\\
48.8667 & 36.8490548018 & 1.39113450 & MS\\
\bottomrule\end{tabular}\end{center}
\textbf{Strong-coupling fits} ($\sigma=s^{-1/2}$; coefficients of $1,\sigma,\sigma^2,\dots$ for
free fits, of $\sigma,\sigma^2,\dots$ after subtracting $4/3$ for fixed fits):
\begin{center}\small\begin{tabular}{llll}\toprule fit & $s_{\min}$ & terms & coefficients\\\midrule
free & 15 & 4 & 1.33285, 0.46502, -0.30670, -0.67069\\
free & 15 & 5 & 1.33382, 0.44473, -0.14893, -1.20874, 0.67936\\
$a_1=4/3$ & 15 & 3 & 0.45752, -0.26849, -0.73442\\
$a_1=4/3$ & 15 & 4 & 0.45483, -0.22703, -0.94389, 0.34674\\
free & 20 & 4 & 1.33309, 0.46096, -0.28422, -0.71170\\
free & 20 & 5 & 1.33361, 0.44934, -0.18722, -1.06855, 0.48851\\
$a_1=4/3$ & 20 & 3 & 0.45692, -0.26188, -0.75242\\
$a_1=4/3$ & 20 & 4 & 0.45554, -0.23883, -0.87926, 0.23015\\
free & 25 & 4 & 1.33321, 0.45884, -0.27176, -0.73598\\
free & 25 & 5 & 1.33344, 0.45342, -0.22291, -0.93062, 0.28970\\
$a_1=4/3$ & 25 & 3 & 0.45666, -0.25873, -0.76166\\
$a_1=4/3$ & 25 & 4 & 0.45591, -0.24528, -0.84159, 0.15735\\
\bottomrule\end{tabular}\end{center}
\textbf{Global approximant} (\texttt{strong/global2.py}). With $t=(1+s/\kappa)^{-1/2}$,
\[ m(s)\approx\frac{3s}{4}+\sum_{k=1}^{10}c_k\,t^k , \]
which is analytic at $s=0$ and expands in $s^{-1/2}$ at infinity. Imposed exactly: $m(0)=3/2$,
$m'(0)=1$, $c_1\sqrt\kappa=4/3$, $c_2\kappa=0.456527$; the other six $c_k$ and
$\kappa=0.3562$ are fitted. Maximum relative error 1.1e-05 on
$0<s\le49$. Coefficients: 2.233910, 1.281505, -0.714086, 6.058693, -75.205525, 230.189503, -352.315510, 303.962417, -142.241360, 28.250452.
It is an interpolant, not a closed form; it does not reproduce the complex singularity at
$|s_*|=0.055$, so its Taylor coefficients beyond $s^1$ are not those of $m(s)$.


\section*{5. What is and is not established}
\textbf{Exact:} $d_5,d_6$ (conditional on the zeta identifications); the algebraic throat and
$g_H(s)$; the location of $s_*$; $a_1=4/3$ \emph{given} the BTZ dictionary for the rotating core.
\textbf{Numerical:} $m(s)$ for $0<s\le\SmaxNum$ with the error budget of \S4;
$a_1=4/3$ to $4\times10^{-4}$ (free fits); $a_2\in[\AtwoLo,\AtwoHi]$ with $a_1$ fixed, consistent with the
conditional analytic value $0.456527$ (supported, not derived here); $a_3\in[\AthreeLo,\AthreeFixHi]$.
\textbf{Conjectural:} that $s_*$ is the nearest singularity of $m(s)$. \textbf{Not obtained:} a closed form for $m(s)$.

\vfill
{\footnotesize Code: \texttt{extremal-mass/code} (exact series), \texttt{extremal-mass/series},
\texttt{extremal-mass/strong}; resume notes \texttt{extremal-mass/RESUME.txt}.}
\end{document}
